% Fuente: Pauta del Control 2, Pregunta 3, ítem 5.
{\color{MidnightBlue}
\begin{enumerate}
    \item \textbf{Falsa.} La convexidad, diferenciabilidad y unicidad del mínimo no garantizan la convergencia del descenso de gradiente para cualquier tamaño de paso fijo $\alpha>0$.

    Considérese la función convexa y diferenciable
    \[
        f(x)=x^2,
        \qquad x\in\mathbb{R}.
    \]
    Esta función tiene un único mínimo global en
    \[
        x^*=0.
    \]
    El método de descenso de gradiente con tamaño de paso fijo $\alpha>0$ genera la sucesión
    \[
        x_{k+1}=x_k-\alpha \nabla f(x_k).
    \]
    Como $\nabla f(x)=2x$, se obtiene
    \[
        x_{k+1}=(1-2\alpha)x_k.
    \]
    Si se toma, por ejemplo, $\alpha=2$ y $x_0\neq 0$, entonces
    \[
        x_{k+1}=-3x_k,
    \]
    de donde
    \[
        x_k=(-3)^k x_0.
    \]
    Esta sucesión no converge a $0$; de hecho, su valor absoluto crece sin límite. Por tanto, no es cierto que el algoritmo converja siempre a la solución óptima para todo $\alpha>0$.
\end{enumerate}
}
